\section{Ce que les trente-cinq trades racontent}

\lettrine{L}{e} moteur or est le contributeur dominant du portefeuille~: il pèse
10\,\% de l'allocation et produit près de 80\,\% du résultat net. Une telle
concentration mérite qu'on regarde les transactions une par une plutôt que de
s'arrêter aux agrégats. C'est l'objet de ce document.

Le matériau est le journal d'exécution du backtest MT5 publié --- le même run
que celui du rapport d'investissement, vérifié par empreinte du fichier de
transactions. Trente-cinq positions y figurent, du 9~novembre~2021 au
21~avril~2026, avec pour chacune le prix d'entrée, le prix de sortie, le volume
en lots, le résultat, le portage, ainsi que le score de décision et le levier
relevés dans le journal du simulateur. Tous les chiffres de ce document
proviennent de cette source unique.

\subsection*{Cinq constats}

\begin{enumerate}[leftmargin=1.4em, itemsep=0.45em]

\item \textbf{Le moteur perd deux fois sur trois et gagne quand même.} Onze
trades gagnants sur trente-cinq, soit 31,4\,\%. Ce n'est pas un défaut~: le gain
moyen ($+$3\,933\,\$) vaut \metric{8,5~fois} la perte moyenne ($-$463\,\$). Un
momentum se paie en petites pertes fréquentes contre de rares grands gains, et
c'est exactement la forme observée.

\item \textbf{Cinq trades font tout le résultat.} Les cinq positions tenues plus
de cent jours rapportent $+$41\,096\,\$, soit \metric{128\,\%} du résultat net
total. Les vingt-et-une positions de moins d'une semaine perdent ensemble
$-$6\,188\,\$. Le moteur ne gagne pas en tradant~: il gagne en \emph{tenant}.

\item \textbf{Le portage coûte 28\,\% du résultat brut.} Le gain de prix
s'élève à $+$44\,926\,\$~; le portage en retire $-$12\,784\,\$. C'est le prix
mécanique de la tenue longue, et il tombe intégralement sur les positions
longues --- celles qui rapportent.

\item \textbf{Il faut encaisser dix pertes d'affilée.} La plus longue série
perdante court sur 255~jours, soit plus de huit mois sans un seul gain. La plus
longue série gagnante, elle, ne compte que deux trades. Tenir la stratégie
pendant ces périodes est une condition d'accès au rendement, pas un détail de
tempérament.

\item \textbf{Le ciblage de volatilité ne cible jamais sa cible.} La volatilité
réalisée de l'or oscille entre 8 et 33\,\% sur la période, toujours sous
l'objectif de 55\,\%. Le levier reste donc structurellement haut --- médiane
5,2$\times$, plafond de 6,6$\times$ atteint onze fois. C'est ce qui produit le
rendement annoncé, et c'est aussi ce qui produit le repli.

\end{enumerate}

\begin{warningbox}
Ces trente-cinq trades proviennent d'un backtest, pas d'un historique réel, et
le simulateur a tourné en mode \emph{OHLC~1~minute}, qui interpole les
exécutions et flatte les résultats. Trente-cinq observations ne permettent
aucune conclusion statistique par sous-groupe. Le chapitre~\ref{sec:limites}
détaille ce que ce document ne démontre pas.
\end{warningbox}

% Généré par scripts/build_gold_trades_figures.py — NE PAS ÉDITER À LA MAIN.
\begin{table}[H]
\centering
\begin{tabular}{@{}lr@{}}
\toprule
\textbf{Métrique} & \textbf{Backtest MT5} \\
\midrule
Trades & 35 \\
Gagnants & 11 (31.4\,\%) \\
Résultat net & +32,141\,\$ \\
Résultat de prix & +44,926\,\$ \\
Swap payé & -12,784\,\$ \\
Gain moyen & +3,933\,\$ \\
Perte moyenne & -463\,\$ \\
Meilleur trade & +21,652\,\$ \\
Pire trade & -1,955\,\$ \\
Rendement de prix moyen & +2.14\,\% \\
Durée médiane & 5\,j \\
Durée maximale & 214\,j \\
Repli maximal de balance & 46.1\,\% \\
Stops de sécurité & 1 \\
\bottomrule
\end{tabular}
\caption{La sleeve or en chiffres, telle que le simulateur d'exécution l'a jouée sur la période 2021--2026.}
\end{table}

